%\section[appx]{Additional related work}
%\label{appx:related_work}
%
%The hierarchical paradigm for planning and executing long-term tasks has been extensively explored over the past several decades \citep{BercherAL2019, PateriaSTQ2021}. The central idea is to decompose complex, long-horizon decision-making problems into a hierarchy of simpler subtasks, where a high-level policy operates over abstract actions, selecting subtasks rather than primitive controls, to achieve the overall objective. Each subtask can itself be formulated as a reinforcement learning problem, solved by a lower-level policy that learns to accomplish it \citep{Hengst2010}. The interaction between these hierarchical policies collectively governs the agent's overall behaviour.
%
%By representing long-horizon problems in terms of temporally extended subtasks, hierarchical reinforcement learning and planning effectively shortens the decision horizon, a principle known as temporal abstraction \citep{BartoM2003, Dietterich2000, SuttonPS2000}. Temporal abstraction enables more efficient credit assignment across extended timescales \citep{VezhnevetsOSHJSK2017}, while decomposing tasks into subtasks typically simplifies learning and promotes more structured exploration during training \citep{NachumTLGLL2019}. These advantages make HRL a powerful framework for scaling reinforcement learning to long-horizon domains \citep{DayanH1992, VezhnevetsOSHJSK2017, NachumGLN2018}. Empirical evidence shows that HRL consistently outperforms flat RL methods across diverse settings, including continuous control \citep{FlorensaDA2017, LevyPS2018}, strategic games \citep{VezhnevetsOSHJSK2017, NachumGLN2018}, and robotic manipulation \citep{FoxKSG2017, GuptaKLLH2019}. Notably, several studies attribute these improvements primarily to enhanced exploration facilitated by subgoal-based hierarchical structures \citep{JongHST2008, NachumTLGLL2019}.
%
%
%Given that large language models (LLMs) encode rich semantic knowledge about the world, such knowledge can be highly valuable for guiding embodied agents to perform high-level reasoning and planning. Motivated by this insight, \citet{AhnBBCC2022} introduced a hierarchical planning framework for robot control, in which a low-level agent executes concrete motor actions, such as opening and closing the gripper, while a high-level, LLM-based agent issues abstract commands like ``put the Coke on the counter.'' Building on this idea, subsequent studies have extensively explored the use of LLMs as task planners. These approaches can be broadly categorized into two types: comprehensive \citep{TangYTZFH2023, DalalCCS2024}, where all subgoals are planned in advance, and incremental, where subgoals are generated adaptively, allowing the planner to make dynamic adjustments based on feedback or environmental changes \citep{IchterBCKFHP2023, ZhangHL2023, WangCMLJLM2024}.
%
%
%%Due to the large number of trainable parameters in LLM-based agents, optimizing them by finetuning their weights is difficult (\red{cite}). As a workaround, most of the existing methods choose to improve the LLM agents' behaviour by updating the prompting with feedback from the environments and/or let the LLM propose a list of possible high-level guidance and train a separate value function to pick them \citep{AhnBBCC2022, WangCaiCLJLM2023}.   
%
%Due to the large number of trainable parameters in LLM-based agents, directly optimizing them through fine-tuning is challenging (\red{cite}). As an alternative, most existing approaches enhance the behaviour of LLM agents by refining their prompts based on feedback from the environment, and/or by allowing the LLM to generate a set of high-level action candidates from which a separate value function selects the most appropriate one \citep{AhnBBCC2022, WangCaiCLJLM2023}. In addition to this popular method, \citet{ZhouHZZL2024} instead propose to use an LLM agent as a teacher to guide a small network to propose high-level guidance. In particular, initially, the agent is trained to follow the guidance of an LLM closely, but the training evolves, 
%
%
%
%
%%most of the existing works do not update its parameters 
%
%
%\red{LLM4Teach } \citep{ZhouHZZL2024}
%
%
%
%Open world
%
%Start from the knowledge from LLM prior 
%
%
%Afterwards, fine-tuning using the pure RL way. 
%
%
%
%
%The hierarchical approach to planning and executing long-term tasks has been extensively studied for decades \citep{BercherAL2019, PateriaSTQ2021}. 
%
%
%
%
%\citep{PateriaSTQ2021}
%
%
%
%%Why we use LLM ---> Generalization capability. 
%
%
%LLM may not be necessary after it provides a good starting point. 
%
%
%
%
%Related work we need to differentiate our work with.
%
%
%\citep{WangCMLJLM2024} --- OmniJARVIS: Unified Vision-Language-Action Tokenization Enables Open-World Instruction Following Agents
%
%\citep{WangCaiCLJLM2023} --- Describe, explain, plan and select: interactive planning with large language models enables open-world multi-task agents
%
%\citep{IchterBCKFHP2023} --- Do As I Can, Not As I Say: Grounding Language in Robotic Affordances
%
%
%\citep{ZhaiBLPSTXZM2024} --- Fine-tuning Large Vision-Language Models as Decision-Making Agents via Reinforcement Learning
%
%
%
%
%
%
%
%
%
%\textbf{Text-based Games (Pure Text I/O) with Hierarchical Control}
%
%\begin{itemize}
%    \item \textbf{Generalization in Text-based Games via Hierarchical RL (H-KGA)} --- two-level setup: a high-level meta-policy selects textual subgoals; a low-level goal-conditioned policy executes. \citep{XuFCDZ2021}.
%    \item \textbf{Deep RL with Stacked Hierarchical Attention for Text-based Games (SHA-KG)} --- hierarchical attention over a knowledge graph for reasoning in text games; \citep{XuFCDZZ2020}
%    \item \textbf{Survey on Text-based Game Benchmarks (Jericho, TextWorld, etc.)} --- provides a map of benchmarks and agent architectures for hierarchical control. \citep{OsborneNF2022}
%\end{itemize}
%
%\textbf{LLM Agents with Hierarchical Planners and Low-level Controllers}
%
%\begin{itemize}
%    \item \textbf{EPO: Environment-aware Planning and Operation for Hierarchical LLM Agents} --- explicit two-layer design: subgoal prediction (planner) and low-level action generation; \citep{ZhaoFSK2024}.
%    \item \textbf{HiPlan: Hierarchical Planning for LLM Agents} --- planner adaptively mixes global (long-range) and local guidance; \citep{LiCYL2025}.
%    \item \textbf{HiAgent: Hierarchical Working-Memory Management for LLM Agents} --- uses subgoals to structure memory and control, achieving strong gains on long-horizon textual tasks; \citep{HuCMMPL2025}.
%    \item \textbf{HLA: Hierarchical Language Agent for Real-time Collaboration} --- high-level reasoning and low-level actuator in text loops for human-AI collaboration; \citep{LiuYGXWLW2024}.
%\end{itemize}
%
%\textbf{Embodied Agents with Text-based Planners (Blueprints for Hierarchical Text RL)}
%
%\begin{itemize}
%    \item \textbf{Voyager: An Open-ended Embodied Agent in Minecraft} --- LLM planner decomposes tasks, grows a reusable skill library, and calls learned skills; \citep{WangXJMXZFA2024}.
%%    \item \textbf{SayCan / SayPlan: Language-driven Robot Control} --- high-level text planner with affordance/value grounding for low-level execution; \citep{IrpanHTZBIFDPFH2022}.
%\end{itemize}






\newpage


%\section{Connection to LLM training paradigms}
%\label{appx:cem_llm_connection}
%The Cross-Entropy Method (CEM) shares a close conceptual resemblance to the training paradigm we employ for large language models, which combines rollout generation, verifier-based evaluation, and supervised fine-tuning (SFT, \citealt{DongXGZPCSZ2023}). In CEM, a population of candidate trajectories is sampled, evaluated, and the top-performing (``elite'') samples are used to update the distribution parameters via cross-entropy minimization. Similarly, in the LLM setting, rollouts from the current LLM are evaluated by a verifier or reward model, and the model is fine-tuned on high-quality responses to increase their likelihood. Both frameworks iteratively shift probability mass toward high-reward regions of the response space. Unlike CEM, however, LLM training does not involve explicit planning or sequential decision-making over trajectories. Nevertheless, this analogy highlights the potential of CEM-style optimization as a scalable principle for aligning and enhancing large language models' long-term planning capabilities through iterative sampling and selection.





\section[appx]{Model architectures}
\label{appx:model_arch}
\begin{figure}[!ht]
    \centering
    
    % Global variable for box width (easy adjustment)
    \newcommand{\boxwidth}{0.5\linewidth}
    
    \begin{tikzpicture}[
        every node/.style={font=\footnotesize},
        box/.style={
            draw,
            rounded corners,
            thick,
            align=left,
            inner sep=6pt,
            text width=\boxwidth,
            anchor=center
        },
        arrow/.style={
            -{Latex[length=2mm,width=1.2mm]},
            thick
        },
        node distance=1.1cm
    ]
    
    % Define colors (soft pastel tones)
    \definecolor{encoderColor}{RGB}{220,235,255} % light blue
    \definecolor{policyColor}{RGB}{230,255,230}  % light green
    \definecolor{decoderColor}{RGB}{255,240,220} % light orange
    
    % Encoder block
    \node[box, fill=encoderColor] (encoder) {
        \textbf{Text Encoder}\\
        Input: token sequence $x_{1:T} \in \Nb^T$\\
        Encodes to latent state $z_s \in \Rb^d$
    };
    
    % Policy block
    \node[box, below=of encoder, fill=policyColor] (policy) {
        \textbf{Policy Network}\\
        Input: $z_s \in \Rb^d$\\
        Samples latent $y \sim \mathcal{N}\Big(\mu(z_s), \diag{\sigmav(z_s)^2}\Big)$\\
        where $\mu(\cdot)$ and $\sigmav(\cdot)$ are neural networks. \\
        Outputs compact action embedding $y \in \Nb^d$
    };
    
    % Decoder block
    \node[box, below=of policy, fill=decoderColor] (decoder) {
        \textbf{Text Decoder}\\
        Input: $y \in \Rb^d$\\
        Decodes to token sequence $\hat{x}_{1:T} \in \Nb^T$
    };
    
    % Arrows
    \draw[arrow] (encoder.south) -- node[right]{state latent $z_s$} (policy.north);
    \draw[arrow] (policy.south) -- node[right]{action embedding $y$} (decoder.north);
    
    \end{tikzpicture}
    
    \caption{
        The Network Architecture of the employ and manager agents.
    }
    \label{fig:text-policy-architecture}
\end{figure}
Although the employee and manager agents have different input–output formats, we convert all signals into text to enable unified processing under the full text-based input–output setting. This design allows both agents to share a universal variational sequence-to-sequence architecture. Examples of output are provided in Appx.~\ref{appx:sample_output}.

As illustrated in \cref{fig:text-policy-architecture}, all networks are implemented using LSTMs \citep{HochreiterS1997}. The LSTM-based text encoder compresses the textual world-state description $x_{1:T} \!\in\! \Nb^T$ into a latent representation $z_s \!\in\! \Rb^d$. Conditioned on $z_s$, the policy network produces an action embedding $y \!\in\! \Rb^d$ via a stochastic latent variable
\[
z_a \sim \mathcal{N}\!\big(\mu(z_s), \operatorname{diag}(\sigmav(z_s)^2)\big),
\]
where $\mu(\cdot)$ and $\sigmav(\cdot)$ are neural networks, and sampling is performed using the reparameterization trick \citep{KingmaW13}. The text decoder then conditions on $y$ to generate the next action sequence $\hat{x}_{1:T} \!\in\! \Nb^T$ in natural language form.





%As we adopt a full-text input and output setting, although employee and manager agents are expected to have different formats of input output pairs, we convert them into text format for subsequential processing. In this way, we can use universal variational sequence-to-sequence model to implement both agents. The input and output samplers are included in Appx~\ref{appx:sample_output}.
%
%In particular, all the networks are implemented based on LSTM \citep{HochreiterS1997}. The Text VAE encoder compresses a textual world-state description $x_{1:T} \!\in\! \Nb^T$ into a scalar latent state $z_s \!\in\! \Rb^d$. The policy network consumes $z_s$ and produces an action embedding $y \!\in\! \Rb^d$ through a stochastic latent 
%\[z_a \sim \mathcal{N}\Big(\mu(z_s), \diag{\sigmav(z_s)^2}\Big)\]
%where $\mu(\cdot)$ and $\sigmav(\cdot)$ are neural networks. And the sampling is performed by re-parametrization trick \citep{KingmaW13}. The Text VAE decoder conditions on $y$ to generate the next natural-language action $\hat{x}_{1:T} \!\in\! \Nb^T$.
\section[appx]{Hyperparameters}
\label{appx:hyperparameter}
\subsection{Autoencoder}
We train an LSTM autoencoder to obtain compact 64\text{-}dimensional latent representations of textual state-action sequences in TextCraft. Both the encoder and decoder operate over 256\text{-}dimensional token embeddings and employ a single-layer LSTM with a 512\text{-}dimensional hidden state. The model is optimized for a maximum of 2000 epochs using Adam \citep{KingmaB2015} with a learning rate of $10^{-3}$ and a batch size of 128. During training, reconstructed sequences are truncated to a maximum length of 30 tokens.

\subsection{World Model}
We employ a Transformer encoder operating over 64-dimensional latent states produced by an LSTM autoencoder. The model consists of a 2-layer Transformer with a hidden size of 256, eight attention heads, 1024-dimensional feedforward blocks, and a dropout rate of 0.1, with a \texttt{[CLS]} token used for sequence-level aggregation. Latent encodings are obtained via a single-layer LSTM autoencoder with 256-dimensional embeddings, a 512-dimensional hidden state, and a 64-dimensional bottleneck, trained using teacher forcing. Optimization is performed using Adam \citep{KingmaB2015} with a learning rate of $3\times10^{-4}$ and weight decay of $1\times10^{-5}$, together with a step-decay scheduler that halves the learning rate every 50 epochs. The world model is trained for 200~epochs with a batch size of 64, conditioning on two past states as temporal context.

\subsection{Policy Agent}
Both policy agents take as input the 64-dimensional latent states derived from the pretrained LSTM autoencoder and use a unidirectional LSTM with a 512-dimensional hidden state, 0.1 dropout. The policy is initialized for RL training from a pre-trained checkpoint and trained for a maximum of 2000 epochs with a batch size of 256 for the employee agent, and 16 for the manager agent, rolling out sequences of length 8 and using EMA smoothing ($\tau = 0.1$) together with 10\% random exploration. Training incorporates a replay buffer of capacity $10{,}000$ with reciprocal transition weighting, random replacement, and mixed sampling consisting of 90\% replayed transitions and 10\% fresh rollouts, with replay sampling enabled after $4096$ collected transitions. We retain 15\% of failed trajectories to improve robustness and perform rollouts every 10 update steps. Optimization uses Adam \citep{KingmaB2015} with a learning rate of $1\times 10^{-4}$. Action execution permits one retry per failed step, and all sequences condition on two past states as temporal context.


\newpage

\setlist[enumerate]{itemsep=1pt, topsep=2pt, parsep=0pt}

\newpage
\section[appx]{Sample outputs}
\label{appx:sample_output}
\paragraph{Trajectory 1.}
\textbf{Final goal:} \textcolor{purple}{red wool}

\textit{Available commands:}
\begin{enumerate}[leftmargin=2em]
    \item craft 1 red wool using 1 red dye, 1 white wool
    \item craft 4 magenta dye using 1 blue dye, 2 red dye, 1 white dye
    \item craft 1 loom using 2 planks, 2 string
    \item craft 1 red dye using 1 poppy
    \item craft 1 white wool using 4 string
    \item craft 2 red dye using 1 rose bush
    \item craft 8 red terracotta using 8 terracotta, 1 red dye
    \item craft 1 bow using 3 stick, 3 string
    \item craft 1 crossbow using 2 string, 3 stick, 1 iron ingot, 1 tripwire hook
    \item craft 1 red dye using 1 red tulip
    \item craft 1 green wool using 1 green dye, 1 white wool
    \item craft 1 white bed using 3 white wool, 3 planks
    \item craft 1 red dye using 1 beetroot
    \item craft 1 black wool using 1 black dye, 1 white wool
    \item craft 2 lead using 4 string, 1 slime ball
    \item craft 1 lime wool using 1 lime dye, 1 white wool
\end{enumerate}

\textbf{Policy inference:}
\begin{enumerate}[label=\textbf{Manager:}, leftmargin=2.4em, labelsep=.5em, align=left]
    \item \textcolor{blue}{Goal:} red dye (1)
    \begin{enumerate}[label=\textbf{Employee:}, leftmargin=2.8em, labelsep=.5em, align=left]
        \item \textcolor{teal}{get} 2 string
        \item \textcolor{teal}{get} 1 beetroot
        \item \textcolor{teal}{craft} 1 red dye using 1 beetroot
    \end{enumerate}

    \item \textcolor{blue}{Goal:} red dye (1), white wool (1)
    \begin{enumerate}[label=\textbf{Employee :}, leftmargin=2.8em, labelsep=.5em, align=left]
        \item \textcolor{teal}{get} 8 string
        \item \textcolor{teal}{craft} 1 white wool using 4 string
    \end{enumerate}

    \item \textcolor{blue}{Goal:} red wool (1)
    \begin{enumerate}[label=\textbf{Employee :}, leftmargin=2.8em, labelsep=.5em, align=left]
        \item \textcolor{teal}{craft} 1 red wool using 1 red dye, 1 white wool
    \end{enumerate}
\end{enumerate}

\paragraph{Trajectory 2.}
\textbf{Final goal:} \textcolor{purple}{spruce slab}

\textit{Available commands:}
\begin{enumerate}[leftmargin=2em]
    \item craft 1 spruce button using 1 spruce planks
    \item craft 4 spruce planks using 1 spruce logs
    \item craft 6 spruce slab using 3 spruce planks
    \item craft 3 spruce sign using 6 spruce planks, 1 stick
    \item craft 2 spruce trapdoor using 6 spruce planks
    \item craft 1 spruce boat using 5 spruce planks
    \item craft 1 spruce pressure plate using 2 spruce planks
    \item craft 1 spruce fence gate using 4 stick, 2 spruce planks
\end{enumerate}

\textbf{Policy inference:}
\begin{enumerate}[label=\textbf{Manager :}, leftmargin=2.4em, labelsep=.5em, align=left]
    \item \textcolor{blue}{Goal:} spruce logs (1)
    \begin{enumerate}[label=\textbf{Employee :}, leftmargin=2.8em, labelsep=.5em, align=left]
        \item \textcolor{teal}{get} 1 spruce logs
    \end{enumerate}

    \item \textcolor{blue}{Goal:} spruce planks (3)
    \begin{enumerate}[label=\textbf{Employee :}, leftmargin=2.8em, labelsep=.5em, align=left]
        \item \textcolor{teal}{craft} 4 spruce planks using 1 spruce logs
    \end{enumerate}

    \item \textcolor{blue}{Goal:} spruce slab (1)
    \begin{enumerate}[label=\textbf{Employee :}, leftmargin=2.8em, labelsep=.5em, align=left]
        \item \textcolor{teal}{craft} 6 spruce slab using 3 spruce planks
    \end{enumerate}
\end{enumerate}

\section[appx]{LLM Prompt and Output for Subgoal Proposal and Verification}
\label{appx:llm_prompt_outputs}
\subsection{Subgoal Decomposition -- \texorpdfstring{$f_\text{dc}$}{f\_dc}}
%\section*{Appendix: Chat History (Exact Verbatim Record)}



\textbf{[User]}
\begin{UserTurn}
Definition
A subgoal is defined as a high-level action or decision that, when performed correctly, significantly advances progress toward successfully completing the overall task. Subgoals may consist of one or more intermediate steps that must be achieved in order to reach the ultimate goal.

Task Description
1. Induce a consistent high-level plan or strategy-a sequence of subgoals-based on the provided set of expert trajectories.
2. Summarize the key steps required to successfully complete the task.
3. The number of subgoals should be smaller than the total number of steps in the trajectory (0 < num_subgoals < num_steps). In other words, not every step is a subgoal.
4. Use consistent methodology in determining these subgoals such that the logic is transferrable.
Output Requirements
- Produce a list of subgoals derived from the given trajectories that collectively lead to the overall goal.
- Each subgoal must be expressed in the exact format:
  the required state needed to achieve the subgoal
- The ultimate goal must appear as the final subgoal in the list.
Please complete this task systematically, so that we can later design a program to automate the process.
Input Format
The provided trajectories will have the following structure:
[goal]
[list of (action, outcome_state) tuples]
The trajectories are below:
0.
Goal: craft birch trapdoor.
craft 6 birch slab using 3 birch planks
get 2 birch logs, {'birch logs': 2}
craft 4 birch planks using 1 birch logs, {'birch planks': 4, 'birch logs': 1}
craft 4 birch planks using 1 birch logs, {'birch planks': 8}
craft 2 birch trapdoor using 6 birch planks, {'birch trapdoor': 2, 'birch planks': 2}
1.
Goal: craft black terracotta.
get 8 terracotta, {'terracotta': 8}
get 1 wither rose, {'terracotta': 8, 'wither rose': 1}
craft 1 black dye using 1 wither rose, {'terracotta': 8, 'black dye': 1}
craft 8 black terracotta using 8 terracotta, 1 black dye, {'black terracotta': 8}
2.
Goal: craft birch button.
get 1 birch logs, {'birch logs': 1}
craft 4 birch planks using 1 birch logs, {'birch planks': 4}
craft 1 birch button using 1 birch planks, {'birch planks': 3, 'birch button': 1}
3.
Goal: craft pink banner.
(get 6 pink tulip ... craft 1 pink banner using 6 pink wool, 1 stick)
4.
Goal: craft crossbow.
(get 2 string ... craft 1 crossbow using 2 string, 3 stick, 1 iron ingot, 1 tripwire hook)
\end{UserTurn}


\textbf{[GPT-5]}
\begin{ModelTurn}
Here’s a consistent, recipe-aware decomposition for each trajectory...

# 0) Goal: craft birch trapdoor
- {'birch logs': 1}
- {'birch planks': 6}
- {'birch trapdoor': 1}

# 1) Goal: craft black terracotta
- {'terracotta': 8}
- {'black dye': 1}
- {'black terracotta': 8}

# 2) Goal: craft birch button
- {'birch logs': 1}
- {'birch planks': 1}
- {'birch button': 1}

# 3) Goal: craft pink banner
- {'white wool': 6}
- {'pink dye': 6}
- {'pink wool': 6}
- {'stick': 1}
- {'pink banner': 1}

# 4) Goal: craft crossbow
- {'string': 2}
- {'stick': 4}
- {'iron ingot': 2}
- {'warped planks': 1}
- {'tripwire hook': 1}
- {'crossbow': 1}
\end{ModelTurn}

\textbf{[User]}
\begin{UserTurn}
Good, now create a Python function to programmatically perform this inference
\end{UserTurn}

\textbf{[GPT-5]}
\begin{ModelTurn}
import re
from collections import defaultdict, deque

CRAFT_RE = re.compile(r"^craft\s+(\d+)\s+([a-z ]+?)\s+using\s+(.+)$")
GET_RE = re.compile(r"^get\s+(\d+)\s+([a-z ]+?)$")

def _parse_ingredients(txt):
    """
    '8 terracotta, 1 black dye' -> list of (qty:int, item:str)
    """
    parts = [p.strip() for p in txt.split(",")]
    out = []
    for p in parts:
        m = re.match(r"^(\d+)\s+([a-z ]+?)$", p.strip())
        if not m:
            # be tolerant of stray punctuation (e.g., trailing periods)
            p = p.strip().rstrip(".,;:")
            m = re.match(r"^(\d+)\s+([a-z ]+?)$", p)
        if m:
            out.append((int(m.group(1)), m.group(2).strip()))
    return out

def _parse_action(line):
    """
    Returns a dict with keys:
      kind: 'craft' or 'get' or None
      out_qty, out_item, ing: for craft
      qty, item: for get
    """
    line = line.strip()
    # cut off trailing inventory dict if present: ", {...}"
    if ", {" in line:
        line = line[:line.index(", {")].strip()

    m = CRAFT_RE.match(line)
    if m:
        out_qty = int(m.group(1))
        out_item = m.group(2).strip()
        ing_txt = m.group(3).strip()
        # tolerate a trailing inventory after a comma already removed; also strip trailing periods
        ing_txt = ing_txt.rstrip(".")
        ing = _parse_ingredients(ing_txt)
        return {"kind": "craft", "out_qty": out_qty, "out_item": out_item, "ing": ing}

    m = GET_RE.match(line)
    if m:
        return {"kind": "get", "qty": int(m.group(1)), "item": m.group(2).strip()}

    return {"kind": None}

def infer_subgoals_from_trajectory(goal_item, action_lines):
    """
    Parameters
    ----------
    goal_item : str
        e.g., 'birch trapdoor', 'black terracotta'
    action_lines : list[str]
        Lines like:
          'get 8 terracotta, {...}'
          'craft 8 black terracotta using 8 terracotta, 1 black dye, {...}'

    Returns
    -------
    subgoals : list[dict]
        Each element is a required inventory state (e.g., {'terracotta': 8})
        Ordered from base → intermediates → final goal.
    """
    # 1) Parse actions into recipes and note which items can be crafted.
    recipes = {}  # product_item -> (out_qty, [(qty, item), ...])
    seen_get = set()
    all_actions = []
    for line in action_lines:
        a = _parse_action(line)
        all_actions.append(a)
        if a["kind"] == "craft":
            recipes[a["out_item"]] = (a["out_qty"], a["ing"])
        elif a["kind"] == "get":
            seen_get.add(a["item"])

    # 2) Find the final craft that produces the goal_item (if present).
    # If absent, we still assume we want at least 1 goal_item.
    if goal_item not in recipes:
        # We can't expand dependencies; just return the final subgoal.
        return [{goal_item: 1}]

    # 3) Recursive expansion to compute ingredient requirements for one goal item.
    # We will compute a dependency graph and a "required counts" per node.
    # For items with recipes, we compute how much of each ingredient is needed
    # to yield the requested quantity (ceil div if needed).
    import math

    # store required quantities for intermediates (not bases)
    required = defaultdict(int)
    # frontier items to expand: (item, qty_needed)
    frontier = deque([(goal_item, 1)])

    # Graph for topo order: edges (ingredient -> product)
    edges = defaultdict(set)
    nodes = set()

    # Track per-item total requirement (for all nodes, including base)
    total_required = defaultdict(int)
    total_required[goal_item] = 1

    while frontier:
        item, need_qty = frontier.popleft()
        nodes.add(item)

        if item in recipes:
            out_qty, ing = recipes[item]
            # crafts needed to produce at least need_qty
            crafts_needed = math.ceil(need_qty / out_qty)
            # accumulate requirements for ingredients
            for q, ing_item in ing:
                req = q * crafts_needed
                total_required[ing_item] += req
                edges[ing_item].add(item)
                nodes.add(ing_item)

                # If ingredient itself is craftable, expand further; else it's base.
                if ing_item in recipes:
                    # push what *this node* needs (not cumulative to avoid overcounting)
                    frontier.append((ing_item, req))

    # 4) Identify base vs intermediates (base = no recipe)
    base_items = [n for n in nodes if n not in recipes and n != goal_item]
    inter_items = [n for n in nodes if n in recipes and n != goal_item]

    # 5) Topological order (Kahn) so that base → intermediates → goal.
    indeg = {n: 0 for n in nodes}
    for u in edges:
        for v in edges[u]:
            indeg[v] += 1

    q = deque([n for n in nodes if indeg[n] == 0])
    topo = []
    while q:
        u = q.popleft()
        topo.append(u)
        for v in edges.get(u, []):
            indeg[v] -= 1
            if indeg[v] == 0:
                q.append(v)

    # 6) Build subgoals:
    #    - For each base item appearing before any of its consumers, include {item: required_count}
    #    - For each intermediate item, include {item: required_count}
    #    - End with {goal_item: 1}
    subgoals = []
    for n in topo:
        if n == goal_item:
            continue
        if n in total_required and total_required[n] > 0:
            subgoals.append({n: total_required[n]})

    # Final goal state: have at least 1 of the goal item.
    subgoals.append({goal_item: 1})
    return subgoals
\end{ModelTurn}

\subsection{Subgoal Completion -- \texorpdfstring{$f_\text{sg}$}{f\_\text{sg}}}
%\section*{Appendix: Full Interaction Log Between User and GPT-5}

\textbf{[User]}
\begin{UserTurn}
Definition
A subgoal is defined as a high-level action or decision that, when performed correctly, significantly advances progress toward successfully completing the overall task. Subgoals may consist of one or more intermediate steps that must be achieved in order to reach the ultimate goal.

Task Description
1. Induce a consistent high-level plan or strategy--a sequence of subgoals--based on the provided set of expert trajectories.
2. Summarize the key steps required to successfully complete the task.
3. The number of subgoals should be smaller than the total number of steps in the trajectory (0 < num_subgoals < num_steps). In other words, not every step is a subgoal.
4. Use consistent methodology in determining these subgoals such that the logic is transferrable.
Output Requirements
- Produce a list of subgoals derived from the given trajectories that collectively lead to the overall goal.
- Each subgoal must be expressed in the exact format:
  the required state needed to achieve the subgoal
- The ultimate goal must appear as the final subgoal in the list.
- Add a reward function such that, given a current state, returns whether a proposed subgoal has been achieved or not (binary label)
Please complete this task systematically, so that we can later design a program to automate the process.
Input Format
The provided trajectories will have the following structure:
[goal]
[list of (action, outcome_state) tuples]
The trajectories are below:
0.
Goal: craft birch trapdoor.
craft 6 birch slab using 3 birch planks
get 2 birch logs, {'birch logs': 2}
craft 4 birch planks using 1 birch logs, {'birch planks': 4, 'birch logs': 1}
craft 4 birch planks using 1 birch logs, {'birch planks': 8}
craft 2 birch trapdoor using 6 birch planks, {'birch trapdoor': 2, 'birch planks': 2}
1.
Goal: craft black terracotta.
get 8 terracotta, {'terracotta': 8}
get 1 wither rose, {'terracotta': 8, 'wither rose': 1}
craft 1 black dye using 1 wither rose, {'terracotta': 8, 'black dye': 1}
craft 8 black terracotta using 8 terracotta, 1 black dye, {'black terracotta': 8}
2.
Goal: craft birch button.
get 1 birch logs, {'birch logs': 1}
craft 4 birch planks using 1 birch logs, {'birch planks': 4}
craft 1 birch button using 1 birch planks, {'birch planks': 3, 'birch button': 1}
3.
Goal: craft pink banner.
get 6 pink tulip, {'pink tulip': 6}
craft 1 pink dye using 1 pink tulip, {'pink tulip': 5, 'pink dye': 1}
get 3 bone meal, {'pink tulip': 5, 'pink dye': 1, 'bone meal': 3}
craft 4 birch planks using 1 birch logs, {'pink tulip': 5, 'pink dye': 1, 'bone meal': 3, 'birch planks': 4}
craft 1 pink dye using 1 pink tulip, {'pink tulip': 4, 'pink dye': 2, 'bone meal': 3, 'birch planks': 4}
craft 1 pink dye using 1 pink tulip, {'pink tulip': 3, 'pink dye': 3, 'bone meal': 3, 'birch planks': 4}
craft 1 pink dye using 1 pink tulip, {'pink tulip': 2, 'pink dye': 4, 'bone meal': 3, 'birch planks': 4}
craft 1 pink dye using 1 pink tulip, {'pink tulip': 1, 'pink dye': 5, 'bone meal': 3, 'birch planks': 4}
craft 1 pink dye using 1 pink tulip, {'pink tulip': 0, 'pink dye': 6, 'bone meal': 3, 'birch planks': 4}
get 24 string, {'pink dye': 6, 'bone meal': 3, 'birch planks': 4, 'string': 24}
craft 1 white wool using 4 string, {'pink dye': 6, 'bone meal': 3, 'birch planks': 4, 'string': 20, 'white wool': 1}
craft 1 white wool using 4 string, {'pink dye': 6, 'bone meal': 3, 'birch planks': 4, 'string': 16, 'white wool': 2}
craft 1 white wool using 4 string, {'pink dye': 6, 'bone meal': 3, 'birch planks': 4, 'string': 12, 'white wool': 3}
craft 1 white wool using 4 string, {'pink dye': 6, 'bone meal': 3, 'birch planks': 4, 'string': 8, 'white wool': 4}
craft 1 white wool using 4 string, {'pink dye': 6, 'bone meal': 3, 'birch planks': 4, 'string': 4, 'white wool': 5}
craft 1 white wool using 4 string, {'pink dye': 6, 'bone meal': 3, 'birch planks': 4, 'string': 0, 'white wool': 6}
craft 1 pink wool using 1 pink dye, 1 white wool, {'pink dye': 5, 'bone meal': 3, 'birch planks': 4, 'white wool': 5, 'pink wool': 1}
craft 1 pink wool using 1 pink dye, 1 white wool, {'pink dye': 4, 'bone meal': 3, 'birch planks': 4, 'white wool': 4, 'pink wool': 2}
craft 1 pink wool using 1 pink dye, 1 white wool, {'pink dye': 3, 'bone meal': 3, 'birch planks': 4, 'white wool': 3, 'pink wool': 3}
craft 1 pink wool using 1 pink dye, 1 white wool, {'pink dye': 2, 'bone meal': 3, 'birch planks': 4, 'white wool': 2, 'pink wool': 4}
craft 1 pink wool using 1 pink dye, 1 white wool, {'pink dye': 1, 'bone meal': 3, 'birch planks': 4, 'white wool': 1, 'pink wool': 5}
craft 1 pink wool using 1 pink dye, 1 white wool, {'pink dye': 0, 'bone meal': 3, 'birch planks': 4, 'white wool': 0, 'pink wool': 6}
get 2 bamboo, {'bone meal': 3, 'birch planks': 4, 'pink wool': 6, 'bamboo': 2}
craft 1 stick using 2 bamboo, {'bone meal': 3, 'birch planks': 4, 'pink wool': 6, 'bamboo': 0, 'stick': 1}
craft 1 pink banner using 6 pink wool, 1 stick, {'bone meal': 3, 'birch planks': 4, 'pink banner': 1}
4.
Goal: craft crossbow.
get 2 string, {'string': 2}
get 6 bamboo, {'string': 2, 'bamboo': 6}
craft 1 stick using 2 bamboo, {'string': 2, 'bamboo': 4, 'stick': 1}
craft 4 spruce planks using 1 spruce logs, {'string': 2, 'bamboo': 4, 'stick': 1, 'spruce planks': 4}
craft 1 stick using 2 bamboo, {'string': 2, 'bamboo': 2, 'stick': 2, 'spruce planks': 4}
get 4 spruce planks, {'string': 2, 'bamboo': 2, 'stick': 2, 'spruce planks': 8}
craft 1 stick using 2 bamboo, {'string': 2, 'bamboo': 0, 'stick': 3, 'spruce planks': 8}
get 1 iron ingot, {'string': 2, 'bamboo': 0, 'stick': 3, 'spruce planks': 8, 'iron ingot': 1}
get 2 bamboo, {'string': 2, 'bamboo': 2, 'stick': 3, 'spruce planks': 8, 'iron ingot': 1}
craft 1 stick using 2 bamboo, {'string': 2, 'bamboo': 0, 'stick': 4, 'spruce planks': 8, 'iron ingot': 1}
get 1 iron ingot, {'string': 2, 'bamboo': 0, 'stick': 4, 'spruce planks': 8, 'iron ingot': 2}
get 1 warped stems, {'string': 2, 'bamboo': 0, 'stick': 4, 'spruce planks': 8, 'iron ingot': 2, 'warped stems': 1}
craft 4 warped planks using 1 warped stems, {'string': 2, 'bamboo': 0, 'stick': 4, 'spruce planks': 8, 'iron ingot': 2, 'warped planks': 4}
craft 2 tripwire hook using 1 stick, 1 iron ingot, 1 warped planks, {'string': 2, 'bamboo': 0, 'stick': 3, 'spruce planks': 8, 'iron ingot': 1, 'warped planks': 3, 'tripwire hook': 2}
craft 1 crossbow using 2 string, 3 stick, 1 iron ingot, 1 tripwire hook, {'spruce planks': 8, 'warped planks': 3, 'tripwire hook': 1, 'crossbow': 1}
\end{UserTurn}

\textbf{[GPT-5]}
\begin{ModelTurn}
Got it -- here's a clean, transferable method plus the concrete subgoals for each trajectory and a tiny reward checker you can drop into your pipeline.

# Method (used consistently below)
1) Read the final recipe and list its direct ingredients.  
2) For any ingredient that is itself crafted, compress its internal chain into a single milestone state (e.g., "have 6 pink wool"), rather than exposing every micro-step.  
3) Keep only necessary inventory thresholds that enable the next craft.  
4) Ensure 0 < #subgoals < #steps and end with the ultimate goal state (>=1 of target item).

---

## 0) Goal: craft birch trapdoor
**Key steps (summary):** convert birch logs -> planks; once you have >=6 birch planks, craft trapdoor.

**Subgoals (states only, final one is the goal):**
1. {'birch planks': 6}  
2. {'birch trapdoor': 1}

---

## 1) Goal: craft black terracotta
**Key steps (summary):** obtain terracotta (8); obtain a black dye (e.g., from wither rose); craft black terracotta.

**Subgoals:**
1. {'terracotta': 8}  
2. {'black dye': 1}  
3. {'black terracotta': 1}

---

## 2) Goal: craft birch button
**Key steps (summary):** make birch planks; craft button.

**Subgoals:**
1. {'birch planks': 1}  
2. {'birch button': 1}

---

## 3) Goal: craft pink banner
**Key steps (summary):** produce 6 pink wool (via white wool + pink dye), get a stick, craft banner.

**Subgoals:**
1. {'pink wool': 6}  
2. {'stick': 1}  
3. {'pink banner': 1}

*(Note: This compresses the internal chains "string -> white wool" and "tulip -> pink dye -> pink wool" into the single milestone `pink wool: 6`.)*

---

## 4) Goal: craft crossbow
**Key steps (summary):** craft a tripwire hook; ensure remaining ingredients (2 string, 3 sticks, 1 iron ingot) are on hand; craft crossbow.

**Subgoals:**
1. {'tripwire hook': 1}  
2. {'string': 2, 'stick': 3, 'iron ingot': 1}  
3. {'crossbow': 1}

*(Note: The hook's own chain (stick + iron ingot + planks) is compressed into the milestone ``have 1 tripwire hook''.)*

---

# Binary Reward Function (subgoal achieved?)
from typing import Dict

def subgoal_achieved(current_state: Dict[str, int], subgoal_state: Dict[str, int]) -> int:
    """
    Returns 1 if current_state satisfies (>=) every item requirement in subgoal_state; else 0.
    Both dicts map item_name -> nonnegative int quantity.
    """
    for item, req_qty in subgoal_state.items():
        if current_state.get(item, 0) < req_qty:
            return 0
    return 1

### Example usage
# Example for trajectory 3 (pink banner)
current = {'pink wool': 6, 'stick': 0}
sg1 = {'pink wool': 6}
sg2 = {'stick': 1}
sg3 = {'pink banner': 1}

assert subgoal_achieved(current, sg1) == 1
assert subgoal_achieved(current, sg2) == 0
assert subgoal_achieved({'pink banner': 1}, sg3) == 1

If you want, I can package these into a JSON schema (trajectories -> subgoals array) or wire the checker into your rollout collector.
\end{ModelTurn}

\newpage

%\begin{figure}[!t]
\centering

% --------- TOP: TRAJECTORY BOX ---------
\begin{tcolorbox}[
    colback=blue!2,
    colframe=blue!40,
    colbacktitle=blue!10,
    coltitle=black,
    title=Provided Demonstration Trajectory,
    boxrule=0.3pt,
    sharp corners,
    enhanced jigsaw,
    fonttitle=\bfseries
]
\ttfamily\raggedright
{\bfseries Goal:} craft birch trapdoor.\\[0.2em]
{\bfseries Trajectory (action, state):}\\[0.2em]

craft 6 birch slab using 3 birch planks, \char`\{\char`\}\\
get 2 birch logs, \char`\{`birch logs': 2\char`\}\\
craft 4 birch planks using 1 birch logs, \char`\{`birch planks': 4, `birch logs': 1\char`\}\\
craft 4 birch planks using 1 birch logs, \char`\{`birch planks': 8\char`\}\\
craft 2 birch trapdoor using 6 birch planks, \char`\{`birch trapdoor': 2, `birch planks': 2\char`\}
\end{tcolorbox}
% --------- BOTTOM: TWO EQUAL-HEIGHT PANELS ---------
\begin{tcbraster}[
  raster equal height,
  raster columns=2,
  raster width=\textwidth,
  raster column skip=3mm,
  boxrule=0.3pt,
  colback=blue!2,
  colframe=blue!40,
  colbacktitle=blue!10,
  coltitle=black,
  enhanced jigsaw,
  sharp corners,
  fonttitle=\bfseries
]


% left box = 45% width
\begin{tcolorbox}[title=LLM-Generated Subgoals, width=0.45\textwidth]
\ttfamily\raggedright
\char`\{`birch logs': 1\char`\}\\
\char`\{`birch planks': 6\char`\}\\
\char`\{`birch trapdoor': 1\char`\}
\end{tcolorbox}
\begin{tcolorbox}[title=Hand-Engineered Subgoals, width=0.53\textwidth]
\ttfamily\raggedright
\char`\{`birch planks': 4, `birch logs': 1\char`\}\\
\char`\{`birch planks': 8\char`\}\\
\char`\{`birch trapdoor':\;\;2, `birch planks':\;\;2\char`\}
\end{tcolorbox}
\end{tcbraster}

\caption{Comparison of LLM-generated and hand-engineered subgoal decompositions for the demonstration trajectory shown above. Additional samples are provided in \cref{appx:llmgenVSEngSubgoal}.}
\label{fig:subgoal_seq}
\end{figure}

\subsection{Comparison of LLM-generated and Hand-engineered Subgoal Decompositions}
\label{appx:llmgenVSEngSubgoal}
\begin{comment}
\subsubsection{Trajectory 0:\;\;Crafting Birch Trapdoor}

\textbf{Suboptimal Trajectory:} 

Goal:\;\;craft birch trapdoor. \\
craft 6 birch slab using 3 birch planks \\
get 2 birch logs \\
craft 4 birch planks using 1 birch logs \\
craft 4 birch planks using 1 birch logs \\
craft 2 birch trapdoor using 6 birch planks \\

\begin{center}
\begin{tabular}{p{0.45\linewidth} p{0.45\linewidth}}
\textbf{LLM-Generated Subgoals} & \textbf{Hand-Engineered Subgoals} \\
\begin{verbatim}
{"birch logs":\;\;1}
{"birch planks":\;\;6}
{"birch trapdoor":\;\;1}
\end{verbatim}
&
\begin{verbatim}
{"birch planks":\;\;4}
{"birch planks":\;\;6}
{"birch trapdoor":\;\;2}
\end{verbatim}
\end{tabular}
\end{center}
\end{comment}

\subsubsection{Trajectory 1: Crafting Black Terracotta}

\begin{tcolorbox}[
    colback=blue!2,
    colframe=blue!40,
    colbacktitle=blue!10,
    coltitle=black,
    title=Provided Demonstration Trajectory,
    boxrule=0.3pt,
    sharp corners,
    enhanced jigsaw,
    fonttitle=\bfseries
]
\ttfamily\raggedright
{\bfseries Goal:} \;\;craft black terracotta.\\[0.2em]
{\bfseries Trajectory (action, state):}\\[0.2em]
get 8 terracotta, \char`\{`terracotta':\;\;8\char`\}\\
get 1 wither rose, \char`\{`terracotta':\;\;8, `wither rose':\;\;1\char`\}\\
craft 1 black dye using 1 wither rose, \char`\{`terracotta':\;\;8, `black dye':\;\;1\char`\}\\
craft 8 black terracotta using 8 terracotta, 1 black dye, \char`\{`black terracotta':\;\;8\char`\}
\end{tcolorbox}

\vspace{0.5em}

\begin{tcbraster}[
  raster equal height,
  raster columns=2,
  raster width=\textwidth,
  raster column skip=3mm,
  boxrule=0.3pt,
  colback=blue!2,
  colframe=blue!40,
  colbacktitle=blue!10,
  coltitle=black,
  enhanced jigsaw,
  sharp corners,
  fonttitle=\bfseries
]

\begin{tcolorbox}[title=LLM-Generated Subgoals]
\ttfamily\raggedright
\char`\{`terracotta':\;\;8\char`\}\\
\char`\{`black dye':\;\;1\char`\}\\
\char`\{`black terracotta':\;\;8\char`\}
\end{tcolorbox}
\begin{tcolorbox}[title=Hand-Engineered Subgoals]
\ttfamily\raggedright
\char`\{`terracotta':\;\;8, `black dye':\;\;1\char`\}\\
\char`\{`black terracotta':\;\;8\char`\}
\end{tcolorbox}

\end{tcbraster}



\subsubsection{Trajectory 2: Crafting Birch Button}

\begin{tcolorbox}[
    colback=blue!2,
    colframe=blue!40,
    colbacktitle=blue!10,
    coltitle=black,
    title=Provided Demonstration Trajectory,
    boxrule=0.3pt,
    sharp corners,
    enhanced jigsaw,
    fonttitle=\bfseries
]
\ttfamily\raggedright
{\bfseries Goal:} \;\;craft birch button.\\
{\bfseries Trajectory (action, state):}\\[0.2em]
get 1 birch logs, \char`\{`birch logs':\;\;1\char`\}\\
craft 4 birch planks using 1 birch logs, \char`\{`birch planks':\;\;4\char`\}\\
craft 1 birch button using 1 birch planks, \char`\{`birch planks':\;\;3, `birch button':\;\;1\char`\}
\end{tcolorbox}

\vspace{0.5em}

\begin{tcbraster}[
  raster equal height,
  raster columns=2,
  raster width=\textwidth,
  raster column skip=3mm,
  boxrule=0.3pt,
  colback=blue!2,
  colframe=blue!40,
  colbacktitle=blue!10,
  coltitle=black,
  enhanced jigsaw,
  sharp corners,
  fonttitle=\bfseries
]

\begin{tcolorbox}[title=LLM-Generated Subgoals]
\ttfamily\raggedright
\char`\{`birch logs':\;\;1\char`\}\\
\char`\{`birch planks':\;\;1\char`\}\\
\char`\{`birch button':\;\;1\char`\}
\end{tcolorbox}
\begin{tcolorbox}[title=Hand-Engineered Subgoals]
\ttfamily\raggedright
\char`\{`birch planks':\;\;4\char`\}\\
\char`\{`birch button':\;\;1, `birch planks':\;\;3\char`\}
\end{tcolorbox}

\end{tcbraster}



\subsubsection{Trajectory 3: Crafting Crossbow}

\begin{tcolorbox}[
    colback=blue!2,
    colframe=blue!40,
    colbacktitle=blue!10,
    coltitle=black,
    title=Provided Demonstration Trajectory,
    boxrule=0.3pt,
    sharp corners,
    enhanced jigsaw,
    fonttitle=\bfseries
]
\ttfamily\raggedright
{\bfseries Goal:} \;\;craft crossbow.\\[0.2em]
{\bfseries Trajectory (action, state):}\\[0.2em]
get 2 string, \char`\{`string':\;\;2\char`\}\\
get 6 bamboo, \char`\{`string':\;\;2, `bamboo':\;\;6\char`\}\\
craft 1 stick using 2 bamboo, \char`\{`string':\;\;2, `bamboo':\;\;4, `stick':\;\;1\char`\}\\
craft 4 spruce planks using 1 spruce logs, \char`\{`string':\;\;2, `bamboo':\;\;4, `stick':\;\;1\char`\}\\
craft 1 stick using 2 bamboo, \char`\{`string':\;\;2, `bamboo':\;\;2, `stick':\;\;2\char`\}\\
get 4 spruce planks, \char`\{`string':\;\;2, `bamboo':\;\;2, `stick':\;\;2\char`\}\\
craft 1 stick using 2 bamboo, \char`\{`string':\;\;2, `stick':\;\;3\char`\}\\
get 1 iron ingot, \char`\{`string':\;\;2, `stick':\;\;3, `iron ingot':\;\;1\char`\}\\
get 2 bamboo, \char`\{`string':\;\;2, `bamboo':\;\;2, `stick':\;\;3, `iron ingot':\;\;1\char`\}\\
craft 1 stick using 2 bamboo, \char`\{`string':\;\;2, `stick':\;\;4, `iron ingot':\;\;1\char`\}\\
get 1 iron ingot, \char`\{`string':\;\;2, `stick':\;\;4, `iron ingot':\;\;2\char`\}\\
get 1 warped stems, \char`\{`string':\;\;2, `stick':\;\;4, `iron ingot':\;\;2, `warped stems':\;\;1\char`\}\\
craft 4 warped planks using 1 warped stems, \char`\{`string':\;\;2, `stick':\;\;4, `iron ingot':\;\;2, `warped planks':\;\;4\char`\}\\
craft 2 tripwire hook using 1 stick, 1 iron ingot, 1 warped planks, \char`\{`string':\;\;2, `stick':\;\;3, `iron ingot':\;\;1, `warped planks':\;\;3, `tripwire hook':\;\;2\char`\}\\
craft 1 crossbow using 2 string, 3 stick, 1 iron ingot, 1 tripwire hook, \char`\{`warped planks':\;\;3, `tripwire hook':\;\;1, `crossbow':\;\;1\char`\}
\end{tcolorbox}

\vspace{0.5em}

\begin{tcbraster}[
  raster equal height,
  raster columns=2,
  raster width=\textwidth,
  raster column skip=3mm,
  boxrule=0.3pt,
  colback=blue!2,
  colframe=blue!40,
  colbacktitle=blue!10,
  coltitle=black,
  enhanced jigsaw,
  sharp corners,
  fonttitle=\bfseries
]

\begin{tcolorbox}[title=LLM-Generated Subgoals]
\ttfamily\raggedright
\char`\{`string':\;\;2\char`\}\\
\char`\{`stick':\;\;4\char`\}\\
\char`\{`iron ingot':\;\;2\char`\}\\
\char`\{`warped planks':\;\;1\char`\}\\
\char`\{`tripwire hook':\;\;1\char`\}\\
\char`\{`crossbow':\;\;1\char`\}
\end{tcolorbox}
\begin{tcolorbox}[title=Hand-Engineered Subgoals]
\ttfamily\raggedright
\char`\{`string':\;\;2, `bamboo':\;\;4, `stick':\;\;1\char`\}\\
\char`\{`string':\;\;2, `bamboo':\;\;2, `stick':\;\;2\char`\}\\
\char`\{`string':\;\;2,`stick':\;\;3\char`\}\\
\char`\{`string':\;\;2, `iron ingot':\;\;1, `stick':\;\;4\char`\}\\
\char`{`string':\;\;2, `stick':\;\;4, `iron ingot':\;\;2, `warped planks':\;\;4\char`\}\\
\char`\{`string':\;\;2, `stick':\;\;3, `iron ingot':\;\;1, `warped planks':\;\;3, `tripwire hook':\;\;2\char`\}\\
 \char`\{`warped planks':\;\;3, `tripwire hook':\;\;1, `crossbow':\;\;1\char`\}
\end{tcolorbox}

\end{tcbraster}



%\red{Some output snippets from employee and manager agents.}

%\red{A few words about the architecture}




